\documentclass[11pt]{article}
\usepackage[margin=1in]{geometry}
\usepackage{graphicx}
\usepackage{booktabs}
\usepackage{hyperref}
\usepackage{microtype}
\usepackage{natbib}
\usepackage{caption}
% Generated by manuscript_validation.py; do not edit empirical values manually.
\newcommand{\InitialWindows}{2822}
\newcommand{\IntensityWindows}{2379}
\newcommand{\EfficiencyWindows}{2054}
\newcommand{\MatchCount}{64}
\newcommand{\TeamCount}{32}
\newcommand{\EventCount}{234,637}
\newcommand{\FreezeFrameCount}{203,882}
\newcommand{\PressureEventCount}{16,554}
\newcommand{\PressureFreezeFrameCount}{14,504}
\newcommand{\PoissonDispersion}{1.603}
\newcommand{\IntensityLeadingEffect}{0.968}
\newcommand{\IntensityLeadingCoef}{-0.032}
\newcommand{\IntensityLeadingCILow}{-0.103}
\newcommand{\IntensityLeadingCIHigh}{0.038}
\newcommand{\IntensityLeadingHolmP}{0.497}
\newcommand{\IntensityTrailingEffect}{1.145}
\newcommand{\IntensityTrailingCoef}{0.135}
\newcommand{\IntensityTrailingCILow}{0.042}
\newcommand{\IntensityTrailingCIHigh}{0.229}
\newcommand{\IntensityTrailingHolmP}{0.005}
\newcommand{\IntensityLeadingPrediction}{6.992}
\newcommand{\IntensityDrawingPrediction}{7.220}
\newcommand{\IntensityTrailingPrediction}{8.265}
\newcommand{\IntensityTrailingDifference}{14.470}
\newcommand{\IntensityLeadingDifference}{-3.161}
\newcommand{\EfficiencyLeadingEffect}{0.918}
\newcommand{\EfficiencyLeadingCoef}{-0.086}
\newcommand{\EfficiencyLeadingCILow}{-0.231}
\newcommand{\EfficiencyLeadingCIHigh}{0.060}
\newcommand{\EfficiencyLeadingHolmP}{0.497}
\newcommand{\EfficiencyTrailingEffect}{1.260}
\newcommand{\EfficiencyTrailingCoef}{0.231}
\newcommand{\EfficiencyTrailingCILow}{0.080}
\newcommand{\EfficiencyTrailingCIHigh}{0.382}
\newcommand{\EfficiencyTrailingHolmP}{0.005}
\newcommand{\EfficiencyLeadingPrediction}{0.181}
\newcommand{\EfficiencyDrawingPrediction}{0.194}
\newcommand{\EfficiencyTrailingPrediction}{0.232}
\newcommand{\EfficiencyTrailingDifference}{3.815}
\newcommand{\EfficiencyLeadingDifference}{-1.289}
\newcommand{\DescriptiveIntensityLeadingPooled}{7.143}
\newcommand{\DescriptiveIntensityLeadingMean}{7.708}
\newcommand{\DescriptiveIntensityLeadingMedian}{7.059}
\newcommand{\DescriptiveIntensityDrawingPooled}{7.157}
\newcommand{\DescriptiveIntensityDrawingMean}{8.009}
\newcommand{\DescriptiveIntensityDrawingMedian}{7.105}
\newcommand{\DescriptiveIntensityTrailingPooled}{7.404}
\newcommand{\DescriptiveIntensityTrailingMean}{8.980}
\newcommand{\DescriptiveIntensityTrailingMedian}{8.182}
\newcommand{\DescriptiveEfficiencyLeadingPooled}{0.188}
\newcommand{\DescriptiveEfficiencyLeadingMean}{0.202}
\newcommand{\DescriptiveEfficiencyLeadingMedian}{0.167}
\newcommand{\DescriptiveEfficiencyDrawingPooled}{0.197}
\newcommand{\DescriptiveEfficiencyDrawingMean}{0.215}
\newcommand{\DescriptiveEfficiencyDrawingMedian}{0.200}
\newcommand{\DescriptiveEfficiencyTrailingPooled}{0.207}
\newcommand{\DescriptiveEfficiencyTrailingMean}{0.226}
\newcommand{\DescriptiveEfficiencyTrailingMedian}{0.200}

\title{Pressing the Scoreline: Pressing Intensity and Regain Efficiency Across Match States in the 2022 FIFA World Cup}
\author{World Cup Strategy Lab}
\date{Internal draft v1}
\begin{document}
\maketitle
\begin{abstract}
Pressing is multidimensional: teams can apply pressure frequently without converting pressure episodes into controlled regains, and successful regains need not generate attacking value. We separate these constructs in an observational analysis of the 2022 FIFA World Cup. The primary unit is a team--match five-minute window. A Poisson model relates score state to Pressure-event counts using opponent passes as exposure; a grouped-binomial model relates score state to controlled regains within five seconds per Pressure sequence. Both include match time, tournament context, and team fixed effects with match-clustered uncertainty. Relative to drawing, trailing was associated with a higher Pressure-event rate (IRR \IntensityTrailingEffect) and higher sequence-regain odds (OR \EfficiencyTrailingEffect). Leading estimates were smaller and negative relative to drawing (IRR \IntensityLeadingEffect; OR \EfficiencyLeadingEffect). These estimates describe conditional associations in one tournament. They remain sensitive to provider definitions, opportunity denominators, sequence construction, and game-state endogeneity. The analysis provides a reproducible framework that keeps pressing intensity, regain efficiency, and post-regain value analytically distinct.
\end{abstract}
\section{Introduction}

Football tactics unfold under a changing match context. Time remaining, score difference, tournament stage, opponent behavior, and numerical balance jointly shape the choices available to a team. Score state is especially visible: leading, drawing, and trailing create different incentives, yet teams enter those states through the same competitive process that also reflects team strength and match performance \citep{VERIFY_SCORE_STATE_REFERENCE}. Comparisons across states are therefore informative descriptions of tactical behavior, but they are not randomized comparisons.

Pressing is also not one behavior with one natural metric. A team may apply many provider-annotated Pressure actions, organize those actions into fewer or more episodes, regain the ball at different rates, and create different attacking value after a regain. Classic PPDA summarizes eligible opponent passes relative to a documented set of defensive actions and has an inverse interpretation. Provider Pressure events instead record individual applications of pressure. Neither measure alone states how often an episode yields controlled possession. Treating frequency as success obscures a central tactical distinction \citep{VERIFY_PPDA_REFERENCE,VERIFY_PRESSURE_EVENT_REFERENCE}.

This study contributes an integrated operationalization that separates three stages: pressing intensity, pressure-to-regain efficiency, and post-regain attacking value. It defines event-level and sequence-level opportunities, uses raw numerators with explicit exposures, compares fixed five-minute windows with score-state-homogeneous segments, and preserves denominator-threshold and clustering choices in a reproducible robustness grid. The aim is not to identify a mechanism imposed by score state. It is to estimate conditional relationships in an audited tournament dataset and show how interpretation depends on measurement choices.

The 2022 FIFA World Cup offers a high-stakes observational setting containing every match in a compact competition. That completeness is valuable, while heterogeneity across national teams, stages, opponents, and tactical systems limits generalization. We ask:

\begin{enumerate}
\item[\textbf{RQ1.}] How are pressing intensity and pressing height conditionally associated with score state?
\item[\textbf{RQ2.}] How is five-second pressure-sequence regain efficiency conditionally associated with score state?
\item[\textbf{RQ3.}] How does post-regain attacking value vary across score states? RQ3 is secondary or exploratory.
\end{enumerate}

Figure~\ref{fig:design} summarizes the design. The two confirmatory outcomes address frequency and sequence conversion separately; post-regain outcomes provide supporting context rather than a third primary family.

\begin{figure}[ht]
\centering
\includegraphics[width=.9\linewidth]{../../outputs/figures/pressing_score_state/figure_1.pdf}
\caption{Conceptual study design. The analysis distinguishes score-state associations with Pressure-event intensity, five-second sequence-regain efficiency, and secondary post-regain value. The primary unit is the team--match five-minute window; drawing is the reference; primary uncertainty is match-clustered; mixed states, numerical inequality, extra time, and unreliable denominators are excluded.}
\label{fig:design}
\end{figure}

\section{Related Work}

This section is deliberately a structured placeholder review. Every citation key is listed in \texttt{CITATIONS\_TO\_VERIFY.md}; the draft is not literature-complete.

\subsection{Score state and tactical behavior}
Research on game state commonly distinguishes leading, drawing, and trailing and reports associations with possession, risk, tempo, or spatial behavior \citep{VERIFY_SCORE_STATE_REFERENCE}. This work motivates adjustment for time and team identity while also highlighting endogeneity: the processes that generate the score may be related to subsequent tactical observations.

\subsection{Pressing intensity and PPDA}
PPDA is widely used as an intensity proxy, but definitions vary across action mappings and pitch zones \citep{VERIFY_PPDA_REFERENCE}. Its denominator and inverse direction matter: lower classic PPDA means more eligible defensive actions per opponent pass. A provider Pressure-event rate has a different numerator and should not be described as PPDA.

\subsection{Event-based Pressure annotations}
Provider annotations permit analysis of individual Pressure actions and their locations \citep{VERIFY_PRESSURE_EVENT_REFERENCE}. Such events are observed tactical labels rather than continuous measurements of every defender's behavior. Annotation rules and off-ball visibility constrain interpretation.

\subsection{Possession regain and counterpressing}
Counterpressing and regain research motivates temporal windows linking defensive actions to controlled possession \citep{VERIFY_REGAIN_REFERENCE}. Event-level conversion can double-count a shared episode; sequence construction offers an alternative opportunity unit but introduces a clustering-window choice.

\subsection{Game-state endogeneity in football analytics}
Because score state, opponent response, substitutions, fatigue, and tactical adaptation co-evolve, adjusted regressions remain associational \citep{VERIFY_ENDOGENEITY_REFERENCE}. Fixed effects address stable team differences but do not remove unobserved within-match variation.

\subsection{Event data, freeze frames, and reproducibility}
Event-linked freeze frames enrich spatial context, but visible players are not full tracking observations \citep{VERIFY_360_REFERENCE}. Reproducible event-data work requires explicit provider versions, transformations, missingness rules, and regeneration paths \citep{VERIFY_REPRODUCIBILITY_REFERENCE}.

\section{Data}

\subsection{Source and tournament coverage}
The source is StatsBomb Open Data for the 2022 FIFA World Cup, resolved to competition 43 and season 106 at provider commit \texttt{b0bc9f22dd77c206ddedc1d742893b3bbe64baec}. The canonical dataset contains \MatchCount{} matches, \TeamCount{} teams, and \EventCount{} events. The repository stores acquisition support, manifests, checksums, and processed schemas but does not track raw provider JSON. Use of the data remains subject to StatsBomb's Open Data terms and attribution requirements \citep{VERIFY_STATSBOMB_LICENSE_REFERENCE}.

All 64 matches have event and 360 files. There are \FreezeFrameCount{} events linked to 360 records. For provider Pressure events, \PressureFreezeFrameCount{} of \PressureEventCount{} have linked 360 context. These are event-linked visible freeze frames, not continuous tracking: absence from a frame does not imply absence from the pitch or from tactical relevance. Spatial context is therefore exploratory and is not part of either confirmatory outcome.

\subsection{Canonicalization and temporal-spatial rules}
Provider inputs are converted to canonical matches, lineups, events, and 360 schemas with source identifiers and provenance. Coordinates are transformed to a provider-independent pitch and an attacking-team frame. The attacking direction is applied symmetrically by team and period, preserving opponent-relative interpretation. Event time is normalized across periods while retaining source period and timestamp fields.

Penalty-shootout events are excluded from match-state and tactical analyses. Extra time is available in processed data but excluded from the primary window analysis. Own goals are assigned to the benefiting team for score reconstruction while retaining provider event identity; they update score state at the event boundary under the audited right-closed interval convention. Red cards update numerical state without being treated as ordinary tactical events.

\subsection{Analysis sample}
The fixed-window source has \InitialWindows{} team--match observations. Primary exclusions remove extra-time windows, invalid or mixed score states, numerical inequality, and outcome-specific low or missing denominators. Table~1 describes score-state exposure and opportunities; Figure~\ref{fig:coverage} visualizes eligible-window coverage. Null rates remain null when an opportunity denominator is absent and are not converted to zero.

\begin{table}[ht]
\centering
\small
\caption{Sample and descriptive characteristics. Team--match five-minute windows; regular-time, homogeneous score-state and numerical-equality restrictions apply. Values are descriptive and drawing is not used as a model reference in this table.}
\begin{tabular}{llllll}
\hline
State & Common windows & Opponent passes & Pressure events & Sequences & 5-s regains \\
Leading & 632 & 15405 & 3668 & 2972 & 559 \\
Drawing & 1260 & 32605 & 7779 & 6424 & 1268 \\
Trailing & 632 & 14384 & 3550 & 2891 & 599 \\
\hline
\end{tabular}
\end{table}


\begin{figure}[ht]
\centering
\includegraphics[width=.85\linewidth]{../../outputs/figures/pressing_score_state/figure_2.pdf}
\caption{Sample coverage by score state. Counts refer to eligible team--match five-minute windows after regular-time, homogeneous-state, and numerical-equality restrictions. Outcome-specific pass and sequence thresholds yield the model samples reported in Methods. This is descriptive coverage, not adjusted inference.}
\label{fig:coverage}
\end{figure}

\section{Methods}

\subsection{Units, exposure, and exclusions}
The primary analytical unit is team $\times$ match $\times$ fixed five-minute window. Score state is the majority state in the window, restricted to leading, drawing, or trailing, with drawing as reference. Primary windows must be in periods 1 or 2, contain one score state, and begin and end at numerical equality. The intensity outcome additionally requires at least five opponent passes; the efficiency outcome requires at least three Pressure sequences and valid successes not exceeding trials. The principal alternative unit is a score-state-homogeneous segment.

\subsection{Pressing intensity}
The first confirmatory outcome is the raw count of provider Pressure events. A Poisson log-link model uses $\log(\text{opponent passes})$ as an offset. Results are reported as incidence-rate ratios (IRRs), adjusted predictions per 30 opponent passes, and differences from drawing. Classic PPDA remains separate: it is opponent eligible passes divided by mapped defensive actions and excludes provider Pressure events from the classic numerator. Material Poisson overdispersion (\PoissonDispersion) triggered the preregistered NB2 robustness model without replacing the Poisson primary specification.

\subsection{Pressure-to-regain efficiency}
Pressure events for a team are grouped into episodes when consecutive actions are no more than two seconds apart. The second confirmatory outcome counts sequences followed by controlled possession within five seconds. A grouped-binomial logit model uses successful sequence regains as successes and Pressure sequences as trials. Sequence level is primary because several Pressure events can belong to one tactical episode; treating each event as an independent opportunity would change the estimand.

\subsection{Secondary and exploratory outcomes}
Secondary measures include Pressure sequences per opponent possession, classic PPDA components, high-pressure share, mean and median Pressure height, provider counterpress share, event-level five-second regain, three- and eight-second sequence regains, and high-regain rate. Post-regain shot production uses successful five-second regains as its opportunity denominator. Same-possession xG and xT per successful regain are exploratory continuous outcomes; zero-opportunity rows remain missing. Same-possession goal was unavailable in the analysis-ready high-regain table, and provider-augmented PPDA was not promoted to primary or secondary inference.

\subsection{Covariates and uncertainty}
Both primary formulas include score state, centered match minute, its square, competition stage, group matchday, and team fixed effects. Standard errors are clustered by match. Robustness specifications use team clustering, match fixed effects where identifiable, exact goal difference, stage subsets, alternative exposure definitions, mixed-state duration weighting, red-card exclusion, and denominator thresholds. The diagnostics flag one rank deficiency among nuisance indicators; primary contrasts are estimated with a generalized inverse, and the limitation is retained rather than changing the locked formula after estimates were observed.

\subsection{Multiplicity, reliability, and missingness}
Holm adjustment is applied across the two primary outcome families separately for each leading/drawing and trailing/drawing contrast. Effect sizes and intervals remain central; secondary analyses do not inherit primary multiplicity control. Threshold sensitivity evaluates opponent-pass minima of 3, 5, and 10 and sequence minima of 1, 3, and 5 without selecting a threshold by significance. Exclusion flow is generated stepwise, and missing metric values are not silently recoded as zero.

\section{Results}

\subsection{Sample construction}
The source contained \InitialWindows{} team--match five-minute windows from \MatchCount{} matches and \TeamCount{} teams. After the common score-state, period, and numerical-state restrictions and the opponent-pass threshold, \IntensityWindows{} windows entered the intensity model. The sequence denominator restrictions retained \EfficiencyWindows{} windows for the efficiency model. The stepwise exclusion table reconciles every reduction; no match or team was silently removed. Drawing contributed the largest exposure total, while leading and trailing exposure seconds were symmetric in the fixed-window construction. Opportunity totals by state are shown in Table~1.

\subsection{Descriptive results}
Pooled Pressure events per 30 opponent passes were \DescriptiveIntensityLeadingPooled{} when leading, \DescriptiveIntensityDrawingPooled{} when drawing, and \DescriptiveIntensityTrailingPooled{} when trailing. The corresponding unweighted mean window rates were \DescriptiveIntensityLeadingMean{}, \DescriptiveIntensityDrawingMean{}, and \DescriptiveIntensityTrailingMean{}. These differ because pooled rates weight windows by opponent-pass exposure, whereas unweighted summaries give each eligible window equal weight.

Pooled five-second sequence-regain proportions were \DescriptiveEfficiencyLeadingPooled{}, \DescriptiveEfficiencyDrawingPooled{}, and \DescriptiveEfficiencyTrailingPooled{} for leading, drawing, and trailing. Unweighted mean window proportions were \DescriptiveEfficiencyLeadingMean{}, \DescriptiveEfficiencyDrawingMean{}, and \DescriptiveEfficiencyTrailingMean{}. These are unadjusted summaries and should not be read as model-standardized comparisons.

\subsection{Primary pressing-intensity model}
Relative to drawing, trailing was associated with a higher adjusted Pressure-event rate: IRR \IntensityTrailingEffect{} (log-rate coefficient \IntensityTrailingCoef{}, 95\% CI \IntensityTrailingCILow{} to \IntensityTrailingCIHigh{}). Predicted Pressure events per 30 opponent passes were \IntensityTrailingPrediction{} when trailing and \IntensityDrawingPrediction{} when drawing, an adjusted difference of approximately \IntensityTrailingDifference\%. Leading versus drawing had IRR \IntensityLeadingEffect{} (coefficient \IntensityLeadingCoef{}, 95\% CI \IntensityLeadingCILow{} to \IntensityLeadingCIHigh{}); predicted rates were \IntensityLeadingPrediction{} and \IntensityDrawingPrediction{}, respectively, an adjusted difference of approximately \IntensityLeadingDifference\%.

\begin{figure}[ht]
\centering
\includegraphics[width=.82\linewidth]{../../outputs/figures/pressing_score_state/figure_3.pdf}
\caption{Adjusted pressing intensity. Predicted provider Pressure events per 30 opponent passes from the confirmatory Poisson model, using team--match five-minute windows and drawing as reference. Team fixed effects and match-clustered uncertainty are used; mixed state, numerical inequality, extra time, and windows with fewer than five opponent passes are excluded. Exact confidence limits are preserved in the source CSV.}
\label{fig:intensity}
\end{figure}

\subsection{Primary regain-efficiency model}
Trailing versus drawing had an estimated sequence-regain odds ratio of \EfficiencyTrailingEffect{} (log-odds coefficient \EfficiencyTrailingCoef{}, 95\% CI \EfficiencyTrailingCILow{} to \EfficiencyTrailingCIHigh{}). Adjusted predicted probabilities were \EfficiencyTrailingPrediction{} and \EfficiencyDrawingPrediction{}, a \EfficiencyTrailingDifference{} percentage-point difference. Leading versus drawing had OR \EfficiencyLeadingEffect{} (coefficient \EfficiencyLeadingCoef{}, 95\% CI \EfficiencyLeadingCILow{} to \EfficiencyLeadingCIHigh{}); predicted probabilities were \EfficiencyLeadingPrediction{} and \EfficiencyDrawingPrediction{}, a \EfficiencyLeadingDifference{} percentage-point difference.

\begin{figure}[ht]
\centering
\includegraphics[width=.82\linewidth]{../../outputs/figures/pressing_score_state/figure_4.pdf}
\caption{Adjusted five-second Pressure-sequence regain efficiency. Predicted probabilities from the confirmatory grouped-binomial model, using team--match five-minute windows and drawing as reference. Team fixed effects and match-clustered uncertainty are used; mixed state, numerical inequality, extra time, and windows with fewer than three sequences are excluded.}
\label{fig:efficiency}
\end{figure}

\subsection{Multiplicity-adjusted inference}
Holm-adjusted primary-family values were \IntensityTrailingHolmP{} for the trailing intensity contrast and \EfficiencyTrailingHolmP{} for trailing efficiency. The leading values were \IntensityLeadingHolmP{} and \EfficiencyLeadingHolmP{}. Interpretation rests on transformed estimates, adjusted predictions, and interval width in addition to these values.

\begin{table}[ht]
\centering
\small
\caption{Primary adjusted models. Team--match five-minute windows; drawing reference; team fixed effects and match-clustered standard errors; mixed states, numerical inequality, extra time, and unreliable denominators excluded. Confirmatory analyses.}
\begin{tabular}{lllll}
\hline
Outcome & Contrast & Estimate & 95\% CI & Transformed \\
pressure\_events & score\_state[leading] & -0.032 & [-0.103, 0.038] & 0.968 \\
pressure\_events & score\_state[trailing] & 0.135 & [0.042, 0.229] & 1.145 \\
sequence\_regains\_5s/pressure\_sequences & score\_state[leading] & -0.086 & [-0.231, 0.060] & 0.918 \\
sequence\_regains\_5s/pressure\_sequences & score\_state[trailing] & 0.231 & [0.080, 0.382] & 1.260 \\
\hline
\end{tabular}
\end{table}


\subsection{Secondary and exploratory results}
High-pressure share and event-level regain models supplied secondary checks on location and measurement unit. Their estimated trailing contrasts were positive, while the event and sequence estimands remained distinct. Provider counterpress share, high-regain rate, Pressure height, sequence frequency, and classic defensive-action counts are retained in Table~3 rather than selected by statistical appearance. Post-regain shot production is secondary; xG and xT per successful regain are exploratory. Same-possession goal is explicitly unavailable, and provider-augmented PPDA remains an unavailable sensitivity entry in the inferential table.

\begin{table}[ht]
\centering
\small
\caption{Secondary and exploratory adjusted analyses. Team--match five-minute windows; drawing reference; match-clustered uncertainty and outcome-specific exclusions. These analyses are not confirmatory primary outcomes.}
\begin{tabular}{lllll}
\hline
Model & Contrast & Estimate & 95\% CI & Role \\
secondary\_high\_pressure\_share & score\_state[leading] & -0.170 & [-0.326, -0.014] & secondary \\
secondary\_high\_pressure\_share & score\_state[trailing] & 0.266 & [0.117, 0.415] & secondary \\
secondary\_event\_regain\_5s & score\_state[leading] & -0.123 & [-0.278, 0.032] & secondary \\
secondary\_event\_regain\_5s & score\_state[trailing] & 0.215 & [0.064, 0.367] & secondary \\
secondary\_counterpress\_share & score\_state[leading] & -0.163 & [-0.317, -0.010] & secondary \\
secondary\_counterpress\_share & score\_state[trailing] & 0.103 & [-0.054, 0.260] & secondary \\
secondary\_high\_regain\_rate & score\_state[leading] & 0.012 & [-0.226, 0.249] & secondary \\
secondary\_high\_regain\_rate & score\_state[trailing] & 0.104 & [-0.162, 0.370] & secondary \\
secondary\_post\_regain\_shot\_10s & score\_state[leading] & 0.170 & [-0.288, 0.628] & exploratory \\
secondary\_post\_regain\_shot\_10s & score\_state[trailing] & -0.212 & [-0.737, 0.312] & exploratory \\
secondary\_pressure\_sequence\_frequency & score\_state[leading] & 0.080 & [-0.005, 0.164] & secondary \\
secondary\_pressure\_sequence\_frequency & score\_state[trailing] & -0.011 & [-0.097, 0.076] & secondary \\
secondary\_classic\_defensive\_action\_rate & score\_state[leading] & -0.011 & [-0.189, 0.167] & secondary \\
secondary\_classic\_defensive\_action\_rate & score\_state[trailing] & -0.005 & [-0.173, 0.163] & secondary \\
exploratory\_post\_regain\_xg\_per\_success & score\_state[leading] & 0.001 & [-0.004, 0.006] & exploratory \\
exploratory\_post\_regain\_xg\_per\_success & score\_state[trailing] & -0.001 & [-0.007, 0.005] & exploratory \\
exploratory\_post\_regain\_xt\_per\_success & score\_state[leading] & -0.006 & [-0.012, 0.001] & exploratory \\
exploratory\_post\_regain\_xt\_per\_success & score\_state[trailing] & 0.004 & [-0.003, 0.010] & exploratory \\
secondary\_mean\_pressure\_height & score\_state[leading] & -1.285 & [-3.387, 0.817] & secondary \\
secondary\_mean\_pressure\_height & score\_state[trailing] & 4.148 & [2.264, 6.031] & secondary \\
\hline
\end{tabular}
\end{table}


\begin{figure}[ht]
\centering
\includegraphics[width=.82\linewidth]{../../outputs/figures/pressing_score_state/figure_5.pdf}
\caption{Intensity--efficiency comparison. Trailing-minus-drawing adjusted contrasts are shown on interpretable but different scales: percentage change in Pressure events per opponent-pass exposure and percentage-point change in sequence-regain probability. Team--match five-minute windows, drawing reference, primary exclusions, team fixed effects, and match-clustered uncertainty apply. Both outcomes are confirmatory, but their scales are not interchangeable.}
\end{figure}

\section{Robustness Analyses}

The robustness grid preserved all 16 preregistered groups plus the NB2 model triggered by overdispersion. Table~4 records executed and unavailable rows. Agreement is assessed through coefficient direction, transformed magnitude, interval width, sample-size change, convergence, and changes in the measurement target—not through a binary significance rule.

Fixed-window results were compared with state-homogeneous segments using opportunity-appropriate offsets. Mixed windows were also retained under duration-proportional allocation. Separate specifications excluded every red-card-affected observation, restricted the sample to group-stage or knockout matches, and reconfirmed the regular-period restriction. Denominator analyses used opponent-pass thresholds of 3, 5, and 10 and Pressure-sequence thresholds of 1, 3, and 5.

Intensity exposure was alternatively defined using opponent possessions and opponent-possession seconds. These versions address a different opportunity scale and are not expected to reproduce the opponent-pass coefficient mechanically. Regain conversion was redefined at event level and with three- and eight-second sequence windows. These analyses distinguish sensitivity to temporal definition from sensitivity to analytical unit.

Uncertainty was recalculated with team clusters, and match fixed effects were used where identifiable. Exact goal difference replaced the categorical state exposure in a separate robustness model. The NB2 intensity contrast retained the count-model interpretation while allowing variance to exceed the Poisson mean. Several outcome/specification combinations are structurally inapplicable—for example, an intensity-exposure change does not define an alternative efficiency model—and are retained as unavailable rather than omitted.

Full residual diagnostics were not recomputed uniformly for every grid variant. The grid preserves sample size, team and match coverage, clustering, convergence, coefficients, intervals, and explicit diagnostic limitations. This constrains fine-grained influence comparisons across variants and is not treated as evidence of universal model stability.

\begin{table}[ht]
\centering
\small
\caption{Robustness summary. Fixed-window and alternative units/exposures; drawing reference except exact goal difference; clustered uncertainty follows each specification. Executed and unavailable rows are retained.}
\begin{tabular}{lll}
\hline
Specification group & Status & Contrast rows \\
alternative\_minimum\_denominators & executed & 12 \\
event\_level\_regain & executed & 2 \\
event\_level\_regain & unavailable & 1 \\
exact\_goal\_difference & executed & 2 \\
exclude\_extra\_time & executed & 4 \\
exclude\_red\_card\_affected & executed & 4 \\
group\_stage\_only & executed & 4 \\
homogeneous\_state\_segments & executed & 4 \\
knockout\_only & executed & 4 \\
match\_fixed\_effects & executed & 4 \\
mixed\_state\_duration\_proportions & executed & 4 \\
negative\_binomial\_overdispersion & executed & 2 \\
opponent\_possession\_count\_exposure & executed & 2 \\
opponent\_possession\_count\_exposure & unavailable & 1 \\
opponent\_possession\_seconds\_exposure & executed & 2 \\
opponent\_possession\_seconds\_exposure & unavailable & 1 \\
primary\_five\_minute\_windows & executed & 4 \\
sequence\_regain\_3s & executed & 2 \\
sequence\_regain\_3s & unavailable & 1 \\
sequence\_regain\_8s & executed & 2 \\
sequence\_regain\_8s & unavailable & 1 \\
team\_clustered\_uncertainty & executed & 4 \\
\hline
\end{tabular}
\end{table}


\begin{figure}[ht]
\centering
\includegraphics[width=.82\linewidth]{../../outputs/figures/pressing_score_state/figure_6.pdf}
\caption{Robustness comparison. Executed primary and alternative specifications include fixed windows, homogeneous segments, denominator thresholds, alternative exposures, temporal regain windows, clustering choices, match fixed effects, and exact goal difference. Unavailable combinations remain in the source table. Drawing is the categorical reference where applicable; exclusions and uncertainty follow each specification.}
\end{figure}

\section{Discussion}

\subsection{Main interpretation}
In this tournament, trailing states were conditionally associated with more frequent provider Pressure activity per opponent pass and higher estimated odds that a Pressure sequence was followed by controlled possession within five seconds. The two associations coexist, but neither identifies why the pattern appeared. Score state, opponent protection of a lead, field position, substitutions, and opportunity selection evolve together.

\subsection{Intensity and efficiency are distinct}
The simultaneous positive trailing contrasts are analytically meaningful because they arise from different numerators and denominators. Intensity asks how often Pressure events occurred given opponent passing opportunities. Efficiency asks how often a sequence converted given Pressure episodes. Tactical urgency, greater pressing commitment, altered field position, opponents accepting lower-risk possession, and selective initiation of Pressure are plausible hypotheses. They were not directly measured as mechanisms and should be tested with richer designs.

\subsection{Leading states}
Leading estimates were modestly negative relative to drawing for both primary outcomes, with intervals that included the reference-equivalent value on the coefficient scale. This does not establish equivalence or absence of a difference. It indicates that the available tournament sample produced less precise and smaller adjusted contrasts than for trailing.

\subsection{Measurement implications}
Pressure events, Pressure sequences, and classic PPDA answer different questions. Event counts reflect provider annotations; sequences reduce repeated actions into episodes; classic PPDA maps a documented defensive-action set to opponent passes and is inversely scaled. Denominator choice changes both reliability and interpretation. Fixed windows offer comparable time blocks but discard mixed states; homogeneous segments preserve state purity with varying duration. Analysts should state these choices before comparing teams or contexts.

\subsection{Practical relevance}
For analysts and coaches, the framework suggests monitoring activity, conversion, and post-regain value separately. A high activity rate need not imply efficient conversion, and efficient conversion need not imply valuable attacks. The estimates do not prescribe a tactical response. They provide benchmarks for asking whether an observed change reflects opportunity volume, episode conversion, or subsequent attack quality.

\subsection{Future 2026 extension}
A future milestone may connect pressing decisions with dynamic qualification pressure under the 48-team 2026 format. Such work would require prospective qualification-state modeling and safeguards against future-information leakage. No 2026 simulation or empirical result is included here.

\section{Limitations}

The study is observational. Score state is endogenous to prior performance, and tactical behavior can precede and follow scoring. Opponent responses, substitutions, fatigue, and tactical adjustments occur simultaneously, so direction and mechanism are not identified. Unobserved within-match context can remain after adjustment, and team fixed effects remove only stable between-team differences.

The sample covers one international tournament with 32 national teams and 64 matches. National teams differ in preparation time, player pools, tactical systems, and schedule. Clustered uncertainty based on 64 matches may be sensitive to a small number of influential matches, and results may not generalize to domestic leagues, other tournaments, or later provider releases.

Pressure is provider-defined. The two-second sequence rule and five-second controlled-regain rule are analytical definitions, not universal tactical boundaries. Event-derived opponent passes, possessions, and possession time approximate opportunity rather than continuous control. Minimum denominator thresholds trade coverage for reliability, and conclusions can vary under alternative thresholds and exposures.

StatsBomb 360 contains event-linked visible freeze frames rather than continuous tracking. Coverage is incomplete overall and for Pressure events, and off-camera players cannot be inferred. The confirmatory models do not use 360 context; any spatial extension would require coverage analysis and visible-area adjustment.

The model diagnostics identify a one-column rank deficiency among nuisance indicators. A generalized inverse yields reproducible primary contrasts, but individual nuisance coefficients are not uniquely parameterized. Robustness-specific residual and influence diagnostics were not uniformly recomputed, limiting model-by-model diagnostic comparison. Missing values remain missing rather than zeros, but selective missingness may still affect eligible samples.

Finally, post-regain shots, xG, and xT use smaller effective samples and opportunity-normalized definitions. Goal conversion was unavailable in the analysis-ready table. Secondary and exploratory findings are not protected by the Holm procedure used for the two primary families. All conclusions are therefore restricted to adjusted observational associations in this dataset.

\section{Conclusion}

This study separates pressing intensity from pressure-to-regain efficiency and post-regain value in the 2022 FIFA World Cup. Relative to drawing, trailing was associated with higher adjusted Pressure-event frequency per opponent pass and higher estimated five-second sequence-regain odds; leading contrasts were smaller and negative. The results do not identify a score-state mechanism. Their main contribution is a denominator-aware, sequence-aware, reproducible framework that makes measurement and robustness choices explicit. Future work can test these relationships in broader competitions and, separately, examine dynamic qualification pressure without treating the present tournament associations as universal.

\bibliographystyle{plainnat}
\bibliography{references}
\end{document}
